\section{継続物体検出とオープン語彙検出の地形図}
\label{sec:detection_ovd}

本章では，本研究の基盤検出器であるMM-Grounding DINO \cite{mmgdino} を中心に，オープン語彙検出（OVD）の系譜，継続学習と物体検出の交差点，および評価基盤と適応ギャップの三領域を整理する．特に，二段検出器で観測された「忘却は分類側に局在し，回帰側は頑健」という知見 \cite{nsgp,redistill} が，クエリベースかつ視覚-言語対照ヘッドを持つMM-Grounding DINOへ転移するか否かを，本研究で検証すべき仮説として位置づける．


\subsection{オープン語彙検出の定式化と系譜}
\label{subsec:ovd_lineage}

\paragraph{問題設定と動機．}
従来の閉集合検出器は，学習時に固定したカテゴリ集合 $\mathcal{C}_{\mathrm{base}}$ のみを検出できる．これに対しオープン語彙検出は，推論時に任意のテキスト語彙 $\mathcal{C}_{\mathrm{test}}$（$\mathcal{C}_{\mathrm{base}} \subseteq \mathcal{C}_{\mathrm{test}}$ とは限らない）を入力として与え，学習時に箱アノテーションを見ていない新規カテゴリ（novel）も検出することを目指す．動機は，物体検出の語彙を分類器の重み行列から切り離し，視覚-言語事前学習で得た意味空間に接続することで，スケーラブルな汎化を得る点にある．

\paragraph{中核アイデアと系譜．}
OVDのサーベイ \cite{ovdsurvey} に従うと，主要な技術系譜は次のように整理できる．
\begin{enumerate}
\item 視覚-意味マッピング：領域特徴をCLIP \cite{clip} などの埋め込み空間へ射影し，クラス名のテキスト埋め込みとの類似度で分類する（RegionCLIP \cite{regionclip}）．
\item 領域認識のグラウンディング学習：検出を句グラウンディングとして再定式化し，画像-テキスト対の大規模コーパスで領域-語の対応を学習する（GLIP \cite{glip}，Grounding DINO \cite{groundingdino}）．
\item 擬似ラベルと弱教師：画像分類データやキャプションから擬似的な箱ラベルを生成して語彙を拡張する（Detic \cite{detic}）．
\item 蒸留と転移：CLIPの画像-テキスト整合性を検出器へ蒸留する，あるいはViTベースの検出ヘッドを軽量に転移する（OWL-ViT \cite{owlvit}，YOLO-World \cite{yoloworld}）．
\end{enumerate}
これらは相互排他的でなく，現代の基盤検出器はグラウンディング学習と擬似ラベル，蒸留を組み合わせる．

\paragraph{定式化．}
OVDの共通骨格は，領域特徴 $\mathbf{v} \in \mathbb{R}^{d}$ とテキスト埋め込み $\{\mathbf{t}_{c}\}_{c \in \mathcal{C}_{\mathrm{test}}}$ の整合度に基づく分類である．カテゴリ $c$ の予測スコアは内積（対照ロジット）
\begin{equation}
s_{c} = \frac{\langle \mathbf{v},\, \mathbf{t}_{c}\rangle}{\tau},\qquad
p_{c} = \sigma\!\left(s_{c}\right),
\label{eq:ovd_logit}
\end{equation}
で与えられる．ここで $\tau$ は温度，$\sigma$ はシグモイド（または語彙上のソフトマックス）である．閉集合検出器が学習可能な分類器重み $\mathbf{w}_{c}$ を用いるのに対し，OVDでは $\mathbf{t}_{c}$ をテキストエンコーダの出力で置換することで，語彙を推論時に差し替え可能にする点が本質である．

\paragraph{評価設定．}
OVDの代表的評価はOV-COCO／OV-LVIS（base/novelの分割，novelのAP）と，多様なドメインを横断するODinW（Object Detection in the Wild）である．ODinW-13上でGrounding DINOはゼロショットで49.2 mAPを達成する \cite{rf100vl}．

\paragraph{長所と短所．}
長所は語彙の柔軟性とゼロショット汎化である．短所は，事前学習分布から離れた専門ドメイン（医療，顕微鏡，電磁波など）で性能が急落すること（後述 \ref{subsec:rf100}），およびテキスト依存ゆえにクラス名の言語的曖昧性に敏感なことである．

\paragraph{本研究への含意．}
本研究の基盤MM-Grounding DINOはグラウンディング学習系に属し，式 \eqref{eq:ovd_logit} の対照ロジットを検出ヘッドに持つ．ドメイン増加型継続学習において保護すべき資産は，閉集合分類器の重みではなく，テキスト埋め込み空間と領域-語整合の汎化能力（ZCOCO）である点が，閉集合の継続物体検出とは本質的に異なる．


\subsection{基盤検出器MM-Grounding DINOの構造}
\label{subsec:mmgdino}

\paragraph{位置づけと動機．}
MM-Grounding DINO \cite{mmgdino} は，Grounding DINO \cite{groundingdino} をオープンソースの大規模データで再実装・拡張した基盤検出器であり，公開された学習レシピと多ドメイン評価を備える．本研究はこれを継続学習の対象VLMとして採用する．

\paragraph{構造（モジュール分解）．}
MM-Grounding DINOは次の要素から構成される．
\begin{enumerate}
\item 画像バックボーン：Swin Transformer によるマルチスケール視覚特徴抽出．
\item テキストエンコーダ：BERT によるクラス名・プロンプトの埋め込み．
\item 特徴強調器（feature enhancer）：変形可能自己注意（画像側）とテキスト自己注意，および双方向のクロスモーダル注意により，視覚特徴とテキスト特徴を相互に条件づける．
\item 言語誘導クエリ選択（language-guided query selection）：テキストとの類似度が高い画像トークンをデコーダのクエリ初期化に選ぶ機構．
\item クロスモダリティデコーダ：DINO型のクエリを画像・テキスト双方の注意で更新する．
\item 対照ヘッド：各クエリの出力と各テキストトークン埋め込みの内積で，クエリ-語の整合ロジットを算出する（式 \eqref{eq:ovd_logit} に対応）．箱回帰は別ブランチで座標を予測する．
\end{enumerate}

\paragraph{定式化．}
$N$ 個のデコーダ出力クエリ $\{\mathbf{q}_{i}\}_{i=1}^{N}$ と $L$ 個のテキストトークン埋め込み $\{\mathbf{t}_{j}\}_{j=1}^{L}$ に対し，整合ロジット行列 $\mathbf{S} \in \mathbb{R}^{N \times L}$ は
\begin{equation}
S_{ij} = \langle \mathbf{q}_{i},\, \mathbf{t}_{j}\rangle,
\label{eq:mmg_logit}
\end{equation}
で与えられ，カテゴリ $c$ の語が占めるトークン集合 $\mathcal{T}_{c}$ 上の集約（最大または平均）でクエリ $i$ のクラス $c$ スコアを得る．分類損失は焦点損失型の対照損失，箱は $L_{1}$ とGIoUの和で，二部マッチングにより教師と対応づける．

\paragraph{テキストトークン長の制約．}
対照ヘッドは全カテゴリ名を結合したテキスト系列に依存するため，トークン長上限（既定 \texttt{max\_text\_len}$=256$）を超える多クラスドメインでは失敗する．本研究ではconfigで $512$ 等へ拡張する（顕微鏡や実世界の多クラス設定）．

\paragraph{長所と短所．}
長所は，分類器が学習可能重みでなくテキスト埋め込みに置換されているため，語彙差し替えとゼロショットが自然なこと，公開レシピで再現性が高いことである．短所は，モデルが大きく，専門ドメインへの適応にはファインチューニングが必要で，その際に事前学習汎化（ZCOCO）と過去ドメインの忘却が生じる点である．

\paragraph{本研究への含意．}
継続学習の介入対象は，バックボーン・特徴強調器・デコーダ・対照ヘッドのいずれにも置きうる．二段検出器の知見 \cite{nsgp,redistill} を踏まえると，忘却がクエリ-語整合（対照ヘッド・テキスト側）に局在するのか，箱回帰ブランチが頑健なのかを，モジュール局在化分析で確認することが第一の問いとなる（\ref{subsec:nsgp} 参照）．


\subsection{CL-DETR：DETR系継続学習と汎用蒸留・リプレイの限界}
\label{subsec:cldetr}

\paragraph{問題設定と動機．}
増分物体検出（IOD）は，段階ごとに新カテゴリのアノテーションを与えて検出器を更新する設定であり，破滅的忘却を被る \cite{cldetr}．従来は知識蒸留（KD）と模範例リプレイ（ER）で対処してきたが，CL-DETR \cite{cldetr} はこれらを Deformable DETR や UP-DETR などの最新Transformer検出器へ素朴に適用しても十分に機能しないことを指摘した．

\paragraph{中核アイデア．}
CL-DETRは二点の検出器固有適応を導入する．第一に検出器知識蒸留（DKD）で，旧モデルの予測のうち最も情報量が高く信頼できるものに絞って蒸留し，冗長な背景予測を無視し，現タスクの正解ラベルと両立させる．第二に較正済みリプレイで，観測したカテゴリ分布を保つよう模範例をサンプルする．

\paragraph{手順と評価．}
COCO 2017上のIOD設定（例：$40+40$ や $70+10$ のカテゴリ分割）で評価し，当時のSOTAを達成した \cite{cldetr}．

\paragraph{長所と短所．}
長所は，DETR系に蒸留・リプレイを有効化したこと，および「汎用CL手法をそのまま検出器へ適用すると失敗する」ことを明確化した点である．短所は，設定がクラス増分中心であり，純粋なドメイン増分やOVDのゼロショット保持を扱わないこと，模範例保持が前提でプライバシ・記憶制約と相反しうることである．なお，これは2023年時点の知見であり現行SOTAとして引用するものではない．

\paragraph{本研究への含意．}
MM-Grounding DINOもDETR系クエリ検出器であるため，汎用KD/ERの素朴適用が不十分という教訓は本研究に直接適用される．対照ヘッドを持つVLM検出器では，蒸留対象を「旧モデルの整合ロジット分布」とするか「テキスト埋め込み空間」とするかが設計上の分岐点となる．


\subsection{二段IOD忘却の解剖とNSGP（null空間勾配射影）}
\label{subsec:nsgp}

\paragraph{問題設定と動機．}
NSGPを提案する研究 \cite{nsgp} は，二段検出器（Faster R-CNN）における忘却がどのモジュールに局在するかを解剖した．動機は，忘却を一様な現象と仮定する従来手法への反証である．

\paragraph{中核となる知見．}
忘却はRoIヘッドの分類器に支配的に局在し，箱回帰器は段階を超えて頑健である \cite{nsgp}．具体的には，後段で劣化した提案を用いてもタスク $\mathcal{D}_{1}$ のmAPは $\mathbf{P}_{1}$ の $77.4\%$ から $\mathbf{P}_{4}$ の $76.1\%$ へ約 $1.3\%$ しか低下せず（RPN由来の忘却は軽微），高IoU提案（IoU $> 0.7$）を推論時に除去しても結果は $74.5\%$ から $76.4\%$ と同程度で，回帰ブランチの頑健性を示した \cite{nsgp}．

\paragraph{定式化（null空間勾配射影）．}
NSGPは，旧入力が張る部分空間と直交する方向のみで特徴抽出器を更新する．旧段階までの入力（特徴）から非心共分散を逐次累積し，
\begin{equation}
\mathbf{C}_{t} \;=\; \mathbf{C}_{t-1} \;+\; \mathbf{X}_{t}^{\top}\mathbf{X}_{t},
\label{eq:nsgp_cov}
\end{equation}
に対し特異値分解
\begin{equation}
\mathbf{C}_{t} \;=\; \mathbf{U}_{t}\,\Lambda_{t}\,\mathbf{U}_{t}^{\top},
\label{eq:nsgp_svd}
\end{equation}
を行う．最小の特異値に対応する固有ベクトル群 $\mathbf{U}_{t}'$ をnull空間として選び，射影行列
\begin{equation}
\mathcal{B} \;=\; \mathbf{U}_{t}'\,(\mathbf{U}_{t}')^{\top},
\label{eq:nsgp_proj}
\end{equation}
を構成する．勾配 $\mathbf{G}$ は
\begin{equation}
\Delta\mathbf{W} \;=\; \mathbf{G}\,\mathcal{B},
\label{eq:nsgp_update}
\end{equation}
として旧入力部分空間へ干渉しない方向へ射影される．ここで $\mathbf{G}\,\mathcal{B} \perp \mathrm{span}(\mathbf{X}_{\mathrm{old}})$ が（近似的に）成り立ち，旧タスク表現を保つ \cite{nsgp}．

\paragraph{長所と短所．}
長所は，忘却の局在を実証し，分類側の表現を勾配射影で保護する明快な原理を与えたことである．短所は，二段検出器に固有のモジュール分解（RPN／RoI分類器／回帰器）に依拠しており，クエリベース・対照ヘッドの検出器へそのまま移植できる保証がないことである．

\paragraph{本研究への含意．}
式 \eqref{eq:nsgp_cov}--\eqref{eq:nsgp_update} のnull空間射影は，提案手法で検討する勾配部分空間直交化（InfLoRA型）と同族の原理であり，事前学習（COCO）を含む過去ドメインの入力部分空間を保護対象に組み込む発想に直結する．加えて，「忘却は分類側に局在・回帰は頑健」という知見が，MM-Grounding DINOのクエリ表現・対照ヘッド・テキストエンコーダに転移するかは自明でなく，本研究で検証すべき中心的仮説である．転移するならば対照ヘッド・テキスト側を集中的に保護すればよく，転移しないならばクエリ表現全体やクロスモーダル注意の保護が必要となる．


\subsection{蒸留の再検討：RPN蒸留の有効性と分類ヘッドの過信}
\label{subsec:redistill}

\paragraph{問題設定と動機．}
増分検出における知識蒸留は，どの出力を教師として用いるかで効果が大きく変わる \cite{redistill}．この研究は，蒸留対象を出力の種別ごとに再検討した．

\paragraph{中核となる知見．}
RPN（領域提案）の蒸留は有効である一方，分類ヘッドの蒸留は旧モデルが過信した誤った教師信号を与え，かえって性能を損なう \cite{redistill}．対策として，外れ値に頑健な適応的Huber損失を蒸留に用いることで，過信誤教師の影響を抑える．評価はクラス増分に加えてドメイン増分でも有効性が示された \cite{redistill}．

\paragraph{定式化．}
教師信号 $y^{\mathrm{old}}$ と生徒出力 $y$ の差 $r = y - y^{\mathrm{old}}$ に対し，閾値 $\delta$ のHuber損失
\begin{equation}
\mathcal{L}_{\delta}(r) =
\begin{cases}
\tfrac{1}{2}\,r^{2}, & |r| \leq \delta,\\[2pt]
\delta\,\big(|r| - \tfrac{1}{2}\delta\big), & |r| > \delta,
\end{cases}
\label{eq:huber}
\end{equation}
を用い，$\delta$ を蒸留対象の信頼度に応じて適応的に設定することで，過信した分類出力（大きな $|r|$）の寄与を線形に抑える．

\paragraph{長所と短所．}
長所は，「分類ヘッドは忘れやすく，かつ蒸留教師としても信頼できない」という \ref{subsec:nsgp} と整合する知見を蒸留の観点から補強したことである．短所は，二段検出器を前提とし，RPN／分類ヘッドの分離に依存する点である．

\paragraph{本研究への含意．}
分類側の脆弱性が二段検出器で繰り返し確認されることは，VLM検出器の対照ヘッド（クエリ-語整合）を蒸留・正則化の主要対象とする設計動機となる．ただし，過信誤教師の問題は，旧モデルの整合ロジットをそのまま蒸留する素朴な方式の危険性を示唆しており，適応的重みづけ（式 \eqref{eq:huber} 型）の導入余地がある．


\subsection{BPF：増分検出の情報非対称性と段階間整合}
\label{subsec:bpf}

\paragraph{問題設定と動機．}
BPF \cite{bpf} は，増分検出に内在する情報非対称性を問題化する．一枚の画像に過去・現在・未来のカテゴリのインスタンスが共存しうるため，前景／背景の定義が段階ごとに変化し，最適化目的が段階間で不整合になる．

\paragraph{中核アイデア．}
ある段階で背景として学習されたインスタンスが，別段階では前景となるため，素朴な学習は過去・未来クラスを背景として誤って抑圧する．BPFはこの非対称性を補正し，段階間で前景定義を整合させることで，過去クラスの抑圧と未来クラスの取りこぼしを緩和する \cite{bpf}．評価はクラス増分検出ベンチマークで行われる．

\paragraph{長所と短所．}
長所は，増分検出特有の最適化不整合という見落とされがちな問題を定式化したことである．短所は，クラス増分（同一画像に複数段階のクラスが共存）を主眼とし，ドメインが完全に切り替わる本研究の純ドメイン増分とは前提が部分的にずれることである．

\paragraph{本研究への含意．}
ドメイン増分でも，あるドメインの画像に別ドメインのクラスに相当する物体が写りうる（例：実世界画像に文書やスクリーンが写る）ため，前景／背景定義の段階間不整合は完全には無縁でない．OVDでは「背景＝語彙に無い」がテキスト依存で揺れるため，情報非対称性の視点はZCOCO保持の文脈でも検討に値する．


\subsection{TiROD：ロボットのドメインシフトを測るDIL＋CILベンチマーク}
\label{subsec:tirod}

\paragraph{問題設定と動機．}
TiROD \cite{tirod} は，ロボットが遭遇する現実的なドメインシフトを，5つの環境・10タスクで評価するベンチマークである．動機は，実応用ではドメイン増分（DIL）とクラス増分（CIL）が同時に発生する点を反映することである．

\paragraph{中核と評価設定．}
5環境（屋内外など異なる撮影条件）にわたりタスクが進行し，各タスクで新環境・新カテゴリが導入される．DILとCILを同時に課すことで，環境変化への適応と語彙拡張の双方の忘却を測る \cite{tirod}．

\paragraph{長所と短所．}
長所は，ドメイン増分を正面から扱う数少ないベンチマークであり，本研究の設定に最も近いことである．短所は，クラス増分が混在するため純ドメイン増分の分離評価には追加の設計が要ること，および規模が限定的なことである．

\paragraph{本研究への含意．}
本研究の評価設計（ドメイン逐次・忘却・ZCOCO保持の三軸）を構成する際の参照点となる．ただしTiRODは閉集合・小規模であり，VLM検出器のゼロショット保持を測る軸を欠くため，Roboflow100系（\ref{subsec:rf100}）と組み合わせる必要がある．


\subsection{Roboflow100／Roboflow100-VL：多ドメイン適応ギャップの定量化}
\label{subsec:rf100}

\paragraph{問題設定と動機．}
Roboflow100 \cite{rf100} はRoboflow Universe由来の100データセットからなる多ドメイン検出ベンチマークであり，Roboflow100-VL \cite{rf100vl} はこれをVLM／MLLMの汎化評価向けに整備したものである（100データセット，564クラス，7ドメイン）．動機は，インターネット規模の事前学習に乏しい概念（医療・顕微鏡・電磁波など）でVLMがどれだけ失敗するかを定量化することである．

\paragraph{中核となる発見（適応ギャップ）．}
Grounding DINOはODinW-13で49.2 mAPを達成するが，RF100-VLではゼロショットで約16 mAPに留まる \cite{rf100vl}．さらに，挑戦的な医療画像データセットでは，Grounding DINOやQwen2.5-VLのゼロショット精度が2\%未満となり，少数ショットでの概念整合の必要性が示された \cite{rf100vl}．

\paragraph{評価設定と主要数値．}
\begin{table}[h]
\centering
\begin{tabularx}{\linewidth}{Xr}
\hline
評価項目 & 数値 \\
\hline
Grounding DINO ODinW-13 ゼロショット & 49.2 mAP \\
Grounding DINO RF100-VL ゼロショット & 約16 mAP \\
医療画像データセットのゼロショット精度 & $< 2\%$ \\
データセット数／クラス数／ドメイン数 & 100／564／7 \\
\hline
\end{tabularx}
\caption{Roboflow100-VLにおける適応ギャップ \cite{rf100vl}．}
\label{tab:rf100vl}
\end{table}
ODinWからRF100-VLへの約33ポイントの落差が，専門ドメイン適応の必要性（課題A）を定量的に裏づける．

\paragraph{長所と短所．}
長所は，事前学習分布外の専門ドメインを多数含み，VLM検出の適応ギャップを明確に示すことである．短所は，標準では各データセットを独立評価し，ドメイン逐次の継続学習設定や忘却・ZCOCO保持の軸を直接は提供しないことである．

\paragraph{本研究への含意．}
本研究はRoboflow100系の専門ドメイン群をドメイン増分の系列として再構成し，適応（課題A）と忘却（課題B）を同時に測る．適応ギャップが大きいほどファインチューニングによる事前学習汎化の破壊（課題C，ZCOCO低下）も大きくなりうるため，適応と保持のトレードオフを定量化する基盤となる．本プロジェクトのドメイン構成（underwater／aerial／videogames／microscopic／documents／electromagnetic／real\_world）はこのベンチマーク思想に沿う．


\subsection{OVDの頑健性評価：分布シフト下の横並び比較}
\label{subsec:ovdrobust}

\paragraph{問題設定と動機．}
OVD検出器のゼロショット性能は標準ベンチマークで報告されるが，分布シフト下での頑健性は別問題である \cite{ovdrobust}．この研究は，OWL-ViT，YOLO-World，Grounding DINOなどの代表的OVD検出器を，分布シフト（自然な破損・スタイル変化・ドメイン差）下で横並びに評価した．

\paragraph{中核となる知見．}
標準ベンチマークでの強さが分布シフト下の頑健性を必ずしも保証しないこと，検出器・バックボーン・事前学習データによって劣化の様相が異なることが示された \cite{ovdrobust}．

\paragraph{長所と短所．}
長所は，単一指標でなく分布シフトの軸でOVDを評価する視点を与えたことである．短所は，継続学習を扱わず，シフト下の静的評価に留まることである（具体的数値は \cite{ovdrobust} を要確認）．

\paragraph{本研究への含意．}
ゼロショット保持（ZCOCO）の評価では，標準COCOだけでなく分布シフト下での頑健性も併せて見るべきことを示唆する．継続学習による適応が，元来の分布シフト頑健性を損なわないかを測る副次指標として有用である．


\subsection{OWL-ST：Web規模自己学習によるOVDのスケーリング}
\label{subsec:owlst}

\paragraph{問題設定と動機．}
OWL-ST \cite{owlst} は，検出データの希少性を，画像-テキスト対への擬似ラベル付けで克服し，OVDをWeb規模へスケールさせる自己学習レシピである．動機は，画像分類や言語モデルで実現したWeb規模学習を，開世界の局在化に持ち込むことである．

\paragraph{中核アイデアと手順．}
既存検出器で画像-テキスト対に擬似箱を生成し，ラベル空間の選択・擬似アノテーションのフィルタリング・学習効率という三つの課題に対処しながら自己学習する \cite{owlst}．OWLv2アーキテクチャと組み合わせ，10億超の事例へスケールする．

\paragraph{主要数値．}
L/14アーキテクチャにおいて，人間の箱アノテーションを一切見ていないLVIS rareクラスのAPを 31.2 から 44.6 へ（相対43\%）向上させた \cite{owlst}．

\paragraph{長所と短所．}
長所は，擬似ラベル自己学習がOVDの汎化を大きく押し上げることを実証したことである．短所は，膨大な計算資源と大規模画像-テキスト対を要し，専門ドメイン（事前学習に乏しい概念）への効果は別途検証が要ることである．

\paragraph{本研究への含意．}
OWL-STの示す「自己学習による擬似ラベルがゼロショット汎化を底上げする」原理は，継続学習において過去・事前学習ドメインの擬似ラベルやリプレイ代替として活用しうる．本研究では，事前学習汎化（ZCOCO）を保持する正則化と，自己学習的なデータ拡張のどちらが効果的かが設計上の論点となる．


\subsection{本章のまとめと本研究の位置づけ}
\label{subsec:detovd_summary}

OVDは語彙を意味空間へ接続することでゼロショット汎化を獲得したが（\ref{subsec:ovd_lineage}），専門ドメインでの適応ギャップは大きい（\ref{subsec:rf100}）．継続物体検出は，汎用KD/ERの素朴適用が不十分であること（\ref{subsec:cldetr}），忘却が分類側に局在し回帰側が頑健であること（\ref{subsec:nsgp}，\ref{subsec:redistill}），増分特有の情報非対称性（\ref{subsec:bpf}）を明らかにしてきたが，その多くは閉集合・二段検出器・クラス増分を前提とする．本研究は，これらの知見をクエリベースかつ対照ヘッドを持つMM-Grounding DINO（\ref{subsec:mmgdino}）のドメイン増分へ移植し，特に「忘却の分類側局在」がVLM検出器の対照ヘッド・テキスト側へ転移するかを検証する．評価はTiROD（\ref{subsec:tirod}）とRoboflow100系（\ref{subsec:rf100}）の思想を継ぎ，適応（A）・忘却（B）・ZCOCO保持（C）の三軸で行い，null空間勾配射影（式 \eqref{eq:nsgp_proj}）と同族の部分空間直交化を，事前学習（COCO）を保護対象に含めて拡張する点に新規性を置く．

